\chapter{LLM Prompt, Data Availability and Ethics}
\label{app:prompt}
\justifying

The exact prompt and settings of the zero-shot LLM experiment (Section~3.5), from \texttt{src/llm\_baseline.py} and \texttt{configs/config.yaml} in the public repository. One request is sent per batch of 25 comments. The instruction is in English, the comments are in Bangla, and there are no examples. Each \texttt{<comment>} stands for one test comment.

\begin{verbatim}
Classify each Bangla comment as 'toxic' or 'non-toxic'. Reply with
ONLY a numbered list like:
1. toxic
2. non-toxic
...

Comments:
1. <comment 1>
2. <comment 2>
...
25. <comment 25>
\end{verbatim}

\begin{table}[H]
  \centering
  \small
  \begin{tabular}{ll}
    \toprule
    \textbf{Setting} & \textbf{Value} \\
    \midrule
    Provider, model identifier & Groq chat completions API, \texttt{qwen/qwen3.8-27b} \\
    Temperature, reasoning mode & 0, off (\texttt{reasoning\_effort} = \texttt{none}) \\
    Maximum output tokens & 512 (2,048 in the earlier run) \\
    Comments per request, pause & 25, 2 seconds between requests \\
    Attempts per batch & 3, with a pause that grows after each failure \\
    Reply accepted if & exactly one valid label for each comment \\
    Failed batch & left unclassified and excluded from scoring (none in the reported run) \\
    Run dates & 15 November 2026 (reported); 31 October 2026 (earlier, incomplete) \\
    \bottomrule
  \end{tabular}
\end{table}

\section*{Data Availability and Ethics}
\addcontentsline{toc}{section}{Data Availability and Ethics}

The dataset of Belal et al. \cite{belal2023toxic} is public at \url{https://github.com/deepu099cse/Multi-Labeled-Bengali-Toxic-Comments-Classification}. The repository declares no licence, so the data was used for research only and is not redistributed here. The comments were assembled from three earlier datasets. Individual comments are quoted only to explain errors, in paraphrase, with names, handles and links removed. The dataset contains vulgar, hateful and threatening language, which is used only to train and evaluate detection systems, and some quoted examples are offensive. The code, notebooks and instructions are public at \url{https://github.com/mehediparvez/bangla_profanity_word_detection}, with links to the saved checkpoints, and the LLM experiment needs a Groq API key that is not part of the repository. The models make mistakes, as the error analysis shows, and are not meant to remove comments or sanction users without human review.
